\subsection{Structural Diagnosis via Hybrid Knowledge-Based RAG Pipelines}
\label{sec:rag_diagnosis}

\begin{sloppypar}
The final stage of the proposed framework transforms quantitative computer vision observations—specifically crack morphology, orientation, and spatial location—into high-fidelity structural diagnostic reports. This is achieved through a multi-layer Retrieval-Augmented Generation (RAG) agent that utilizes the digitalized knowledge base of **FEMA 306 (Evaluation of Post-Earthquake Damage to Concrete and Masonry Walls)**. The agent bridges the gap between raw pixel-level data and engineering-grade assessment by simulating the reasoning process of a structural expert.
\end{sloppypar}

\subsubsection{Query Formulation and Multi-Modal Input Synthesis}
\begin{sloppypar}
The RAG pipeline is triggered by a structured state vector $\mathcal{V}_{input}$ derived from the preceding vision modules. The input includes:
\end{sloppypar}

\begin{itemize}[noitemsep,leftmargin=1.5em]
    \item \textbf{Morphology}: Crack orientation $\theta^*$ (e.g., Stepped, Diagonal, Vertical) and statistical widths $\{w_{max}, w_{median}\}$.
    \item \textbf{Semantic Location}: The highest-density grid cell $G_{i,j}$ and its intersection with wall components (e.g., Pier, Spandrel).
    \item \textbf{Propagation}: The directional vector $\vec{v}$ indicating the stress origin.
\end{itemize}

\begin{sloppypar}
These parameters are synthesized into a natural language query $\mathcal{Q}$ using a template-based transformation: 
$\mathcal{Q} = \text{"Material: URM, Orientation: } \theta^*, \text{ Location: } G_{i,j}, \text{ Description: crack originating in } P_o \text{ and propagating to } P_t\text{"}$. This query is then embedded into a 384-dimensional dense vector space using the \texttt{all-MiniLM-L6-v2} transformer for semantic retrieval.
\end{sloppypar}

\subsubsection{Two-Layer Reasoning Architecture}
\begin{sloppypar}
To ensure the reliability of the diagnostic output, we implement a structured dual-layer reasoning architecture using the \textbf{Gemini 1.5 Flash} model.
\end{sloppypar}

\paragraph{Layer 1: Failure Mode Identification}
\begin{sloppypar}
The agent retrieves the top-8 most relevant "Classification Guides" from the FEMA 306 knowledge base stored in a ChromaDB vector store. Each candidate represents a discrete failure mechanism (e.g., Bed Joint Sliding, Diagonal Tension, Flexural Rocking). The LLM performs a zero-shot cross-comparison between the field observations and the technical criteria described in each guide. It explicitly filters for morphological alignment (e.g., "X-cracking" maps to Section 7.2.10) and selects the single most probable Failure Mode $\mathcal{M}^*$.
\end{sloppypar}

\paragraph{Layer 2: Severity Assessment and Scoped Reasoning}
\begin{sloppypar}
Once the failure mode is locked, the agent initiates the second reasoning layer. It parses the measured crack width $w$ and maps it to the three-tier severity scale (Fine, Moderate, Severe) defined for the selected mechanism. The agent then performs "Scoped Reasoning," focusing exclusively on the damage level table within the FEMA chunk that matches the calculated severity. This layer generates a detailed explanation of the structural implications—citing specific restoration measures or performance degradation criteria (e.g., "loss of lateral stiffness due to crushing of clay units")—ensuring the output is transparent and auditable by professional engineers.
\end{sloppypar}

\subsubsection{Hybrid Confidence Scoring (HCS)}
\begin{sloppypar}
We propose a Hybrid Confidence Score ($\Phi$) that ensembles three independent reliability metrics to flag potential "hallucinations" or low-confidence assessments. The final score is calculated as a weighted average:
\end{sloppypar}

\begin{equation}
    \Phi = 0.5 \cdot SC_{LLM} + 0.3 \cdot P_{RC} + 0.2 \cdot S_{CE}
\end{equation}

\noindent where:
\begin{itemize}[noitemsep,leftmargin=1.5em]
    \item $SC_{LLM}$ is the Self-Calculated confidence generated by the LLM during Layer 1.
    \item $P_{RC}$ is the Retrieval Confidence, derived from the Softmax normalized L2 distances of the top-8 vector search results.
    \item $S_{CE}$ is the Re-ranked Confidence using a \texttt{ms-marco-MiniLM-L6-v2} Cross-Encoder, which evaluates the semantic similarity between the query and the full text of the retrieved chunks.
\end{itemize}

\begin{sloppypar}
This hybrid approach ensures that the diagnostic report is only generated if the morphological observations are strongly grounded in the engineering knowledge base, providing the "paragraphic" depth required for high-impact assessment.
\end{sloppypar}
